% Escreva aqui a solução da questão 3.
\textbf{VERDADEIRA}

\textcolor{red}{ VERSÃO 1 } %Letícia

%\begin{enumerate}[label=\Roman*)]
%    \item Observe que a diferença entre os conjuntos pode ser reescrita como:
    
%    $\mathcal{U} \setminus A = \mathcal{U} \cap A^{C} = A^{C}$. Pela definição de igualdade de conjuntos, $(\mathcal{U} \setminus A) \subseteq A^{C}$.

%    \item De uma das propriedades da interseção de conjuntos: $(A \cap B) \subseteq A$. 
%    De uma das propriedades do complementar, $(A \cap B) \subseteq A \implies A^{C} \subseteq (A \cap B)^{C}$.

%    Por transitividade, de (I) temos que: $(\mathcal{U} \setminus A) \subseteq (A \cap B)^{C}$. Portanto, a relação de inclusão é verdadeira.
 
%\end{enumerate}
De uma das propriedades da interseção de conjuntos, temos que $(A \cap B) \subseteq A$.
Então, podemos reescrever:
  \begin{align*}
      (A \cap B) \subseteq A &\implies A^{C} \subseteq (A \cap B)^{C} \tag{$X \subseteq Y \implies Y^C \subseteq X^C$} \\
      &\implies (\mathcal{U} \setminus A) \subseteq (A \cap B)^{C}. \tag{$\mathcal{U} \setminus X= X^{C}$}
  \end{align*}

\bigskip

\textcolor{red}{CHAVE DE PONTUAÇÃO:
\begin{itemize}
    \item  Inclusão pela interseção \emph{(1,0 ponto)};
    \item Uso de $X \subseteq Y \implies Y^C \subseteq X^C$ \emph{(1,0 ponto)};
    \item  Uso da igualdade $\mathcal{U} \setminus X= X^{C}$ \emph{(1,0 ponto)}.
\end{itemize} } 

\bigskip




\textcolor{red}{ VERSÃO 2 } %Letícia

Observando o conjunto do lado esquerdo da relação de inclusão, percebe-se que a diferença entre os conjuntos pode ser reescrita como a interseção do conjunto à direita da diferença ($\mathcal{U}$) e do complementar do conjunto à esquerda ($A$). Ou de forma mais direta, a diferença entre o universo e o conjunto $A$ é igual ao complementar de $A$.

\begin{align*}
    \mathcal{U} \setminus A & = \mathcal{U} \cap A^{C} \\
    & = A^{C}.
\end{align*}

No lado esquerdo da inclusão, por De Morgan, temos:
$$(A \cap B)^{C} = A^{C} \cup B^{C}.$$

Então,  $ \mathcal{U} \setminus A \subseteq (A \cap B)^{C}$ é equivalente a $A^{C} \subseteq A^{C} \cup B^{C}$. De uma das propriedades da união, a sentença é verdadeira.

\bigskip



%\textcolor{red}{ VERSÃO 3} %autor





\bigskip

\textcolor{red}{CHAVE DE PONTUAÇÃO:
\begin{itemize}
    \item  Uso da igualdade $\mathcal{U} \setminus X= X^{C}$ \emph{(1,0 ponto)};
    \item  Igualdade pelo De Morgan \emph{(1,0 ponto)};
    \item  Inclusão pela união \emph{(1,0 ponto)}.
\end{itemize} }